% !PRINTED 163
\weberpage{163}

%% [running head, not body text -- mirrored 2026-08-07] Page furniture (running
%% head, folio, signature) is excluded from this edition by standing policy.
%% \noindent\S.\,50.\hfill Elimination from several equations.\hfill 163

%<<CONT
the number of points of intersection of these three surfaces. But if
$T$ is identically zero, then the three surfaces have either a
portion of surface or a curve in common.

We can determine the degree of the final equation in the following
way, which requires no lengthy considerations at all about the
degrees of the functions $R$, $S$, $T$. In order to make ourselves
independent of exceptional cases, we introduce homogeneous variables
by adding a fourth, $u$, and thus consider three homogeneous
quaternary equations
\weq{4}{F(x, y, z, u) = 0, \quad \Phi(x, y, z, u) = 0, \quad \Psi(x, y, z, u) = 0}
of the orders $m$, $n$, $p$.

If we assume the special case in which each of these three functions
splits into purely linear factors, then it is clear that, if the
number of common solutions is not infinite. % [sic] printed period followed by lowercase "their"
their number must be
equal to the product $mnp$; for we can set one linear factor of each
of $F$, $\Phi$, $\Psi$ equal to zero simultaneously, and so obtain
$mnp$ systems of three linear equations each, which, if they are
independent of one another, yield just as many systems of values for
the ratios $x : y : z : u$.

The conclusion that in the general case the degree of the final
equation of (4) is just as great appears questionable at first
sight. But following a remark of \gesperrt{F. Klein}, the doubt can
be completely removed by the following simple consideration.

Let us first imagine the final equation formed from the equations
(4) under the assumption that the coefficients in the equations (4)
are variables independent of one another. Then the resultant which
results after elimination of $x$, $y$ is certainly not identically
zero, since it does not vanish identically even for special systems
of coefficients, e.g.\ when $F$, $\Phi$, $\Psi$ split into linear
factors; and we obtain a homogeneous final equation for the two
unknowns $z$, $u$ of a certain degree, which is to be determined.

If one specialises the coefficients in any way whatever, the final
equation may indeed possibly be satisfied identically; but if this
case does not occur, then its order cannot be changed (since all
terms are of the same order). Now since, by the specialisation
assumed above, which expresses itself in the splitting of the
functions $F$, $\Phi$, $\Psi$, the
%CONT>>
